% GENERATED FILE. Edit canonical lessons, then run npm run generate:books.
% canonical-source-hash: fnv1a64:f733dd2526a83657
% canonical-lessons: TE-C541-natana, TE-C541-sannivesam, TE-C541-kathanayakudu, TE-C541-hasyanatudu, TE-C541-preksakulu, TE-R541-review-534-541

\chapter{Stage and Screen}
\label{ch:te-stage-and-screen}
\hlchaptermodality{541}

\begin{chapteropening}
\textbf{By the end of this chapter:} \emph{I can talk about acting, scenes, heroes, comedians and audiences.}

Complete the last lesson of chapter 541: I can talk about acting, scenes, heroes, comedians and audiences.
\end{chapteropening}

\section[\texorpdfstring{\te{నటన} (naṭana)}{నటన (naṭana)}]{\te{నటన} (naṭana) — acting; a performance}

\label{lesson:TE-C541-natana}



\begin{quote}
Before the new one: say the Telugu for to compose, then the Telugu for a translation.
\end{quote}

\subsection*{You'll want to know: \te{నటన}}
\textbf{\te{నటన}} — \emph{naṭana} — \textquotedblleft{}acting; a performance\textquotedblright{}.

\begin{quote}
\te{ఆ} \te{సినిమాలో} \te{ఆమె} \te{నటన} \te{అద్భుతం}.
\end{quote}

\begin{quote}
\emph{ā sinimālō āme naṭana adbhutaṁ.}
\end{quote}

\begin{quote}
Her acting in that film is wonderful.
\end{quote}

\te{నటుడు} is an actor; \te{నటించు} is to act.

\begin{cousinweb}
\textbf{\te{నటన}} (\emph{naṭana}) is borrowed from Sanskrit \emph{naṭana}, \textquotedblleft{}dancing, dance, pantomime\textquotedblright{} (Monier-Williams).
\end{cousinweb}

\begin{grammarlens}[title={What you've built}]
One more word for describing what you lived and felt.
\end{grammarlens}

\subsection*{Your turn}
\begin{itemize}
\raggedright
  \item \emph{Say it:} \emph{naṭana}
  \item \emph{Say it:} \emph{naṭana}, once more
  \item \emph{Say it:} \emph{ā sinimālō āme naṭana adbhutaṁ.} — Her acting in that film is wonderful.
  \item \emph{Recall:} say the Telugu for to compose, then the Telugu for a translation, then say \emph{naṭana} again
\end{itemize}

\begin{tcolorbox}[breakable,colback=teal!4,colframe=teal!35!black,title={Before you move on}]
What does \te{నటన} mean? (\textquotedblleft{}Acting; a performance\textquotedblright{}.) Say it once more.
\end{tcolorbox}
\section[\texorpdfstring{\te{సన్నివేశం} (sannivēśaṁ)}{సన్నివేశం (sannivēśaṁ)}]{\te{సన్నివేశం} (sannivēśaṁ) — a scene}

\label{lesson:TE-C541-sannivesam}



\begin{quote}
Before the new one: say the Telugu for a translation, then the Telugu for acting.
\end{quote}

\subsection*{You'll want to know: \te{సన్నివేశం}}
\textbf{\te{సన్నివేశం}} — \emph{sannivēśaṁ} — \textquotedblleft{}a scene\textquotedblright{}.

\begin{quote}
\te{సినిమాలో} \te{చివరి} \te{సన్నివేశం} \te{చూసి} \te{అందరూ} \te{ఏడ్చారు}.
\end{quote}

\begin{quote}
\emph{sinimālō civari sannivēśaṁ cūsi andarū ēḍcāru.}
\end{quote}

\begin{quote}
Everyone cried at the last scene of the film.
\end{quote}

A scene in a film or a play; also a situation.

\begin{cousinweb}
\textbf{\te{సన్నివేశం}} (\emph{sannivēśaṁ}) is borrowed from Sanskrit \emph{saṁniveśa}, \textquotedblleft{}settling down together; a position, a situation\textquotedblright{} (Monier-Williams).
\end{cousinweb}

\begin{grammarlens}[title={What you've built}]
One more word for describing what you lived and felt.
\end{grammarlens}

\subsection*{Your turn}
\begin{itemize}
\raggedright
  \item \emph{Say it:} \emph{sannivēśaṁ}
  \item \emph{Say it:} \emph{sannivēśaṁ}, once more
  \item \emph{Say it:} \emph{sinimālō civari sannivēśaṁ cūsi andarū ēḍcāru.} — Everyone cried at the last scene of the film.
  \item \emph{Recall:} say the Telugu for a translation, then the Telugu for acting, then say \emph{sannivēśaṁ} again
\end{itemize}

\begin{tcolorbox}[breakable,colback=teal!4,colframe=teal!35!black,title={Before you move on}]
What does \te{సన్నివేశం} mean? (\textquotedblleft{}A scene\textquotedblright{}.) Say it once more.
\end{tcolorbox}
\section[\texorpdfstring{\te{కథానాయకుడు} (kathānāyakuḍu)}{కథానాయకుడు (kathānāyakuḍu)}]{\te{కథానాయకుడు} (kathānāyakuḍu) — a hero; a male lead}

\label{lesson:TE-C541-kathanayakudu}



\begin{quote}
Before the new one: say the Telugu for acting, then the Telugu for a scene.
\end{quote}

\subsection*{You'll want to know: \te{కథానాయకుడు}}
\textbf{\te{కథానాయకుడు}} — \emph{kathānāyakuḍu} — \textquotedblleft{}a hero; a male lead\textquotedblright{}.

\begin{quote}
\te{ఆ} \te{సినిమాలో} \te{కథానాయకుడు} \te{ఒక} \te{రైతు}.
\end{quote}

\begin{quote}
\emph{ā sinimālō kathānāyakuḍu oka raitu.}
\end{quote}

\begin{quote}
In that film the hero is a farmer.
\end{quote}

Feminine: \te{కథానాయిక}, the heroine.

\begin{cousinweb}
\textbf{\te{కథానాయకుడు}} (\emph{kathānāyakuḍu}) is a learned compound built from Sanskrit \emph{kathā}, \textquotedblleft{}story\textquotedblright{} (CDIAL 2704), and \emph{nāyaka}, \textquotedblleft{}a leader; in drama, the hero\textquotedblright{} (Monier-Williams).
\end{cousinweb}

\begin{grammarlens}[title={What you've built}]
One more word for describing what you lived and felt.
\end{grammarlens}

\subsection*{Your turn}
\begin{itemize}
\raggedright
  \item \emph{Say it:} \emph{kathānāyakuḍu}
  \item \emph{Say it:} \emph{kathānāyakuḍu}, once more
  \item \emph{Say it:} \emph{ā sinimālō kathānāyakuḍu oka raitu.} — In that film the hero is a farmer.
  \item \emph{Recall:} say the Telugu for acting, then the Telugu for a scene, then say \emph{kathānāyakuḍu} again
\end{itemize}

\begin{tcolorbox}[breakable,colback=teal!4,colframe=teal!35!black,title={Before you move on}]
What does \te{కథానాయకుడు} mean? (\textquotedblleft{}A hero; a male lead\textquotedblright{}.) Say it once more.
\end{tcolorbox}
\section[\texorpdfstring{\te{హాస్యనటుడు} (hāsyanaṭuḍu)}{హాస్యనటుడు (hāsyanaṭuḍu)}]{\te{హాస్యనటుడు} (hāsyanaṭuḍu) — a comedian; a comic actor}

\label{lesson:TE-C541-hasyanatudu}



\begin{quote}
Before the new one: say the Telugu for a scene, then the Telugu for a hero.
\end{quote}

\subsection*{You'll want to know: \te{హాస్యనటుడు}}
\textbf{\te{హాస్యనటుడు}} — \emph{hāsyanaṭuḍu} — \textquotedblleft{}a comedian; a comic actor\textquotedblright{}.

\begin{quote}
\te{ఆ} \te{హాస్యనటుడిని} \te{చూస్తేనే} \te{అందరూ} \te{నవ్వుతారు}.
\end{quote}

\begin{quote}
\emph{ā hāsyanaṭuḍini cūstēnē andarū navvutāru.}
\end{quote}

\begin{quote}
The mere sight of that comedian makes everyone laugh.
\end{quote}

\te{హాస్యం}, humour, + \te{నటుడు}, actor.

\begin{cousinweb}
\textbf{\te{హాస్యనటుడు}} (\emph{hāsyanaṭuḍu}) is a compound built from Sanskrit \emph{hāsya}, \textquotedblleft{}laughable; laughter\textquotedblright{} (CDIAL 14071), and \te{నటుడు}, \textquotedblleft{}an actor\textquotedblright{}.
\end{cousinweb}

\begin{grammarlens}[title={What you've built}]
One more word for describing what you lived and felt.
\end{grammarlens}

\subsection*{Your turn}
\begin{itemize}
\raggedright
  \item \emph{Say it:} \emph{hāsyanaṭuḍu}
  \item \emph{Say it:} \emph{hāsyanaṭuḍu}, once more
  \item \emph{Say it:} \emph{ā hāsyanaṭuḍini cūstēnē andarū navvutāru.} — The mere sight of that comedian makes everyone laugh.
  \item \emph{Recall:} say the Telugu for a scene, then the Telugu for a hero, then say \emph{hāsyanaṭuḍu} again
\end{itemize}

\begin{tcolorbox}[breakable,colback=teal!4,colframe=teal!35!black,title={Before you move on}]
What does \te{హాస్యనటుడు} mean? (\textquotedblleft{}A comedian; a comic actor\textquotedblright{}.) Say it once more.
\end{tcolorbox}
\section[\texorpdfstring{\te{ప్రేక్షకులు} (prēkṣakulu)}{ప్రేక్షకులు (prēkṣakulu)}]{\te{ప్రేక్షకులు} (prēkṣakulu) — spectators; viewers}

\label{lesson:TE-C541-preksakulu}



\begin{quote}
Before the new one: say the Telugu for a hero, then the Telugu for a comedian.
\end{quote}

\subsection*{You'll want to know: \te{ప్రేక్షకులు}}
\textbf{\te{ప్రేక్షకులు}} — \emph{prēkṣakulu} — \textquotedblleft{}spectators; viewers\textquotedblright{}.

\begin{quote}
\te{నాటకం} \te{అయ్యాక} \te{ప్రేక్షకులు} \te{లేచి} \te{నిలబడ్డారు}.
\end{quote}

\begin{quote}
\emph{nāṭakaṁ ayyāka prēkṣakulu lēci nilabaḍḍāru.}
\end{quote}

\begin{quote}
When the play ended, the audience stood up.
\end{quote}

Plural; one spectator is \te{ప్రేక్షకుడు}.

\begin{cousinweb}
\textbf{\te{ప్రేక్షకులు}} (\emph{prēkṣakulu}) is borrowed from Sanskrit \emph{prekṣaka}, \textquotedblleft{}looking on; a spectator, a member of an audience\textquotedblright{} (Monier-Williams).
\end{cousinweb}

\begin{grammarlens}[title={What you've built}]
One more word for describing what you lived and felt.
\end{grammarlens}

\subsection*{Your turn}
\begin{itemize}
\raggedright
  \item \emph{Say it:} \emph{prēkṣakulu}
  \item \emph{Say it:} \emph{prēkṣakulu}, once more
  \item \emph{Say it:} \emph{nāṭakaṁ ayyāka prēkṣakulu lēci nilabaḍḍāru.} — When the play ended, the audience stood up.
  \item \emph{Recall:} say the Telugu for a hero, then the Telugu for a comedian, then say \emph{prēkṣakulu} again
\end{itemize}

\begin{tcolorbox}[breakable,colback=teal!4,colframe=teal!35!black,title={Before you move on}]
What does \te{ప్రేక్షకులు} mean? (\textquotedblleft{}Spectators; viewers\textquotedblright{}.) Say it once more.
\end{tcolorbox}
\section[(dialogue)]{Words from chapters 534–541 — a review}

\label{lesson:TE-R541-review-534-541}



\begin{quote}
Say the Telugu for to have a row and a dispute.
\end{quote}

\subsection*{Your turn}
Say these aloud:

\begin{itemize}
\raggedright
  \item to have a row, a dispute, to sulk, a compromise, a rift
  \item mockery, to mock, to reprimand, to call to account, to deceive
  \item a picnic, a pastime, gardening, knitting, a card game
  \item a spinning top, marbles, hide-and-seek, kabaddi, to hop
  \item a musical note, a singing voice, to play an instrument, a mridangam drum, a reed pipe
  \item a refrain, a stanza, a metrical verse, a lyric, folk tradition
  \item literary works, a short story, a piece of writing, to compose, a translation
  \item acting, a scene, a hero, a comedian, spectators
\end{itemize}

\begin{tcolorbox}[breakable,colback=teal!4,colframe=teal!35!black,title={Before you move on}]
To have a row? (\textbf{\te{గొడవపడు}}, \emph{goḍavapaḍu}.) A dispute? (\textbf{\te{తగాదా}}, \emph{tagādā}.) To sulk? (\textbf{\te{అలుగు}}, \emph{alugu}.) A compromise? (\textbf{\te{రాజీ}}, \emph{rājī}.) A rift? (\textbf{\te{మనస్పర్థ}}, \emph{manaspartha}.) Mockery? (\textbf{\te{ఎగతాళి}}, \emph{egatāḷi}.) To mock? (\textbf{\te{వెక్కిరించు}}, \emph{vekkiriñcu}.) To reprimand? (\textbf{\te{మందలించు}}, \emph{mandaliñcu}.) To call to account? (\textbf{\te{నిలదీయు}}, \emph{niladīyu}.) To deceive? (\textbf{\te{మోసగించు}}, \emph{mōsagiñcu}.) A picnic? (\textbf{\te{వనభోజనం}}, \emph{vanabhōjanaṁ}.) A pastime? (\textbf{\te{కాలక్షేపం}}, \emph{kālakṣēpaṁ}.) Gardening? (\textbf{\te{తోటపని}}, \emph{tōṭapani}.) Knitting? (\textbf{\te{అల్లిక}}, \emph{allika}.) A card game? (\textbf{\te{పేకాట}}, \emph{pēkāṭa}.) A spinning top? (\textbf{\te{బొంగరం}}, \emph{boṅgaraṁ}.) Marbles? (\textbf{\te{గోళీలు}}, \emph{gōḷīlu}.) Hide-and-seek? (\textbf{\te{దాగుడుమూతలు}}, \emph{dāguḍumūtalu}.) Kabaddi? (\textbf{\te{చెడుగుడు}}, \emph{ceḍuguḍu}.) To hop? (\textbf{\te{గెంతు}}, \emph{gentu}.) A musical note? (\textbf{\te{స్వరం}}, \emph{svaraṁ}.) A singing voice? (\textbf{\te{గాత్రం}}, \emph{gātraṁ}.) To play an instrument? (\textbf{\te{వాయించు}}, \emph{vāyiñcu}.) A mridangam drum? (\textbf{\te{మృదంగం}}, \emph{mṛdaṅgaṁ}.) A reed pipe? (\textbf{\te{సన్నాయి}}, \emph{sannāyi}.) A refrain? (\textbf{\te{పల్లవి}}, \emph{pallavi}.) A stanza? (\textbf{\te{చరణం}}, \emph{caraṇaṁ}.) A metrical verse? (\textbf{\te{పద్యం}}, \emph{padyaṁ}.) A lyric? (\textbf{\te{గేయం}}, \emph{gēyaṁ}.) Folk tradition? (\textbf{\te{జానపదం}}, \emph{jānapadaṁ}.) Literary works? (\textbf{\te{సాహిత్యం}}, \emph{sāhityaṁ}.) A short story? (\textbf{\te{కథానిక}}, \emph{kathānika}.) A piece of writing? (\textbf{\te{రచన}}, \emph{racana}.) To compose? (\textbf{\te{రచించు}}, \emph{raciñcu}.) A translation? (\textbf{\te{అనువాదం}}, \emph{anuvādaṁ}.) Acting? (\textbf{\te{నటన}}, \emph{naṭana}.) A scene? (\textbf{\te{సన్నివేశం}}, \emph{sannivēśaṁ}.) A hero? (\textbf{\te{కథానాయకుడు}}, \emph{kathānāyakuḍu}.) A comedian? (\textbf{\te{హాస్యనటుడు}}, \emph{hāsyanaṭuḍu}.) Spectators? (\textbf{\te{ప్రేక్షకులు}}, \emph{prēkṣakulu}.)
\end{tcolorbox}
